\documentclass[11pt,a4paper]{article}

\usepackage[brazil]{babel}
\usepackage[T1]{fontenc}
\usepackage{lmodern}
\usepackage{microtype}
\usepackage[margin=2.2cm]{geometry}
\usepackage{amsmath}
\usepackage{booktabs}
\usepackage{tabularx}
\usepackage{enumitem}
\usepackage{xcolor}
\usepackage[hidelinks]{hyperref}

\setlist{nosep, leftmargin=*}
\setlength{\parskip}{0.4em}
\setlength{\parindent}{0pt}

% Trechos que ainda faltam escrever.
\newcommand{\pendente}[1]{\textcolor{red!70!black}{\textbf{[pendente:} #1\textbf{]}}}

\title{\textbf{A* no mapa da Romênia: de Arad a Bucharest}\\[0.3em]
  \large Execução passo a passo}
\author{Luisa Ferreira \quad Caio Trindade\\
  \small Inteligência Artificial para Jogos I --- Aula 26, pathfinding}
\date{Novembro de 2026}

\begin{document}
\maketitle

\section{Problema}

O exercício pede para aplicar o A* no grafo do último slide da aula 26 e calcular o caminho de
\textbf{Arad} até \textbf{Bucharest}. Cada nó $n$ é avaliado por
\[
  f(n) = g(n) + h(n),
\]
em que $g(n)$ é a distância percorrida pelas estradas desde Arad e $h(n)$ é a distância em linha
reta de $n$ até Bucharest, dada pela tabela do slide.

\section{Dados}

\pendente{tabela com as distâncias em linha reta usadas (h) e as estradas relevantes}

\section{Execução}

Convenções: \texttt{OPEN} é a fronteira (nós abertos) e \texttt{CLOSED} são os nós já expandidos.
A cada iteração, o nó de menor $f$ sai de \texttt{OPEN}; se for Bucharest, a busca termina.

\pendente{uma tabela por iteração, com o nó expandido, os filhos gerados (g, h, f, pai) e o
estado de OPEN e CLOSED}

\section{Resultado}

\pendente{caminho encontrado, custo total e por que a primeira entrada de Bucharest na fronteira
não é a solução}

\end{document}
